\chapter{中点与平行}\label{ch:g8:midpoints}

连接三角形两边的中点，所得线段
总与第三边平行，且长度恰为第三边的一半。这
“中点定理”及其逆定理是一个重要思想的初步体验 ---
平行线截得的线段成比例 --- 这一思想将
在\cref{ch:g9:thales}中发展为泰勒斯定理。

\section{中点定理}

\begin{theorem}[中点定理]\label{thm:g8:midpoints:direct}
在三角形 $ABC$ 中，设 $I$ 为 $[AB]$ 的中点，$J$ 为
$[AC]$ 的中点。则直线 $(IJ)$ 平行于 $(BC)$，且
\[
IJ = \frac{BC}{2} .
\]
\end{theorem}

\begin{proof}[证明思路]
设 $K$ 为 $I$ 关于 $J$ 的对称点（即半转，
见\cref{ch:g7:central}）。四边形 $AICK$ 的对角线 $[IK]$
与 $[AC]$ 相交于公共中点 $J$，故它是
平行四边形，于是 $KC$ 平行于 $AI$，即平行于 $(AB)$，且
$KC = AI = IB$。再看 $IBCK$，它有两条对边（$[IB]$ 与
$[KC]$）平行且相等，故它也是平行四边形（由
\cref{thm:g7:central:recognize}可知），因此 $(IK) \parallel (BC)$ ---
即 $(IJ) \parallel (BC)$ --- 且 $BC = IK = 2\,IJ$。
\end{proof}

\begin{omfigure}
\begin{tikzpicture}[scale=1.1]
  \coordinate (A) at (1.7,2.9);
  \coordinate (B) at (0,0);
  \coordinate (C) at (4.6,0);
  \coordinate (I) at ($(A)!0.5!(B)$);
  \coordinate (J) at ($(A)!0.5!(C)$);
  \draw[thick] (A) -- (B) -- (C) -- cycle;
  \draw[omThm, very thick] (I) -- (J);
  \fill (A) circle (1.5pt) node[above, font=\small] {$A$};
  \fill (B) circle (1.5pt) node[below left, font=\small] {$B$};
  \fill (C) circle (1.5pt) node[below right, font=\small] {$C$};
  \fill[omThm] (I) circle (1.5pt) node[left, font=\small] {$I$};
  \fill[omThm] (J) circle (1.5pt) node[right, font=\small] {$J$};
  \draw ($(A)!0.22!(B)$) ++(-0.09,-0.05) -- ++(0.18,0.1);
  \draw ($(A)!0.72!(B)$) ++(-0.09,-0.05) -- ++(0.18,0.1);
  \draw ($(A)!0.25!(C)$) ++(-0.06,-0.1) -- ++(0.12,0.2);
  \draw ($(A)!0.25!(C)$) ++(0.02,-0.1) -- ++(0.12,0.2);
  \draw ($(A)!0.75!(C)$) ++(-0.06,-0.1) -- ++(0.12,0.2);
  \draw ($(A)!0.75!(C)$) ++(0.02,-0.1) -- ++(0.12,0.2);
\end{tikzpicture}

{\small 中点线段 $[IJ]$（红色）平行于第三
边 $(BC)$，且长度为其一半。}
\end{omfigure}

\begin{example}\label{ex:g8:midpoints:direct}
在三角形 $ABC$ 中，$BC = 9$ 厘米，$I$ 为 $[AB]$ 的中点，
$J$ 为 $[AC]$ 的中点，将两者相连。不作任何测量：
则 $(IJ) \parallel (BC)$，且 $IJ = \frac92 = 4.5$ 厘米。小三角形
$AIJ$ 是 $ABC$ 按一半比例缩小的相似形 --- 其周长也是
$ABC$ 周长的一半。
\end{example}

\begin{theorem}[逆定理]\label{thm:g8:midpoints:converse}
在三角形 $ABC$ 中，$I$ 为 $[AB]$ 的中点，过该点
且与 $(BC)$ \emph{平行}的直线与边 $[AC]$ 相交于其中点。
\end{theorem}
\admitted

\begin{example}\label{ex:g8:midpoints:converse}
一条小路穿过一块三角形田地，从某一条边的中点出发，
并与其对边平行。逆定理保证
小路恰好在另一条边的中点处离开田地 --- 另一侧
无需测量。
\end{example}

\begin{method}[在定理与逆定理之间选择]\label{met:g8:midpoints:choose}
\begin{enumerate}
  \item \emph{已知两个中点} $\Rightarrow$ 用直接定理，
        得出平行及长度为一半的结论。
  \item \emph{已知一个中点与一组平行关系} $\Rightarrow$ 用
        逆定理，得出交点是中点。
  \item 写明所用的三角形、中点及定理，
        每项各用一句话 --- 说明。
\end{enumerate}
\end{method}

\section{一半大小的三角形}

\begin{proposition}[中点三角形]\label{prop:g8:midpoints:medial}
连接三角形三边的中点，将其分成
四个面积相等的三角形，每一个都是原三角形一半大小的相似形。
\end{proposition}

\begin{proof}
中点三角形的每条边都平行于 $ABC$ 的某一边，且
长度为其一半（连续三次应用\cref{thm:g8:midpoints:direct}）。
四个小三角形的三边长度相同，故
它们全等；它们共同铺满 $ABC$，因此每个
的面积为原三角形面积的四分之一。
\end{proof}

\begin{omfigure}
\begin{tikzpicture}[scale=1.05]
  \coordinate (A) at (1.9,3.0);
  \coordinate (B) at (0,0);
  \coordinate (C) at (5.0,0);
  \coordinate (I) at ($(A)!0.5!(B)$);
  \coordinate (J) at ($(A)!0.5!(C)$);
  \coordinate (K) at ($(B)!0.5!(C)$);
  \draw[thick] (A) -- (B) -- (C) -- cycle;
  \fill[omThm!12] (I) -- (J) -- (K) -- cycle;
  \draw[omThm, thick] (I) -- (J) -- (K) -- cycle;
  \fill (A) circle (1.4pt) node[above, font=\small] {$A$};
  \fill (B) circle (1.4pt) node[below left, font=\small] {$B$};
  \fill (C) circle (1.4pt) node[below right, font=\small] {$C$};
  \fill[omThm] (I) circle (1.4pt) node[left, font=\small] {$I$};
  \fill[omThm] (J) circle (1.4pt) node[right, font=\small] {$J$};
  \fill[omThm] (K) circle (1.4pt) node[below, font=\small] {$K$};
\end{tikzpicture}

{\small 中点三角形 $IJK$ 将 $ABC$ 分成四个全等的
三角形——一幅值得记住的图。}
\end{omfigure}

\begin{remark}[迈向泰勒斯]
中点定理是一个更一般事实中“一半”的情形：
与三角形一边平行的直线，截另外两边所得的
\emph{相等的比}——三分之一与三分之一，五分之二与
五分之二\,\dots\ 这一般结论即为泰勒斯定理
（见\cref{ch:g9:thales}）；中点的情形是
本学年唯一需要的情形。
\end{remark}

\section{练习}

\begin{exercise}[$\star$]\label{exo:g8:midpoints:1}
在三角形 $RST$ 中，$M$ 为 $[RS]$ 的中点，$N$ 为
$[RT]$ 的中点，已知 $ST = 7$ cm。关于直线
$(MN)$ 与长度 $MN$ 有何结论？写出所用的定理。
\end{exercise}

\begin{exercise}[$\star$]\label{exo:g8:midpoints:2}
在三角形 $ABC$ 中，$I$ 为 $[AB]$ 的中点，
过 $I$ 且平行于 $(BC)$ 的直线交 $[AC]$ 于 $J$，
已知 $AC = 11$ cm。计算 $AJ$，并写出所用的定理。
\end{exercise}

\begin{exercise}[$\star$]\label{exo:g8:midpoints:3}
画三角形 $ABC$，使 $AB = 8$ cm，$AC = 6$ cm，$BC = 7$ cm，
取 $I$ 为 $[AB]$ 的中点，$J$ 为 $[AC]$ 的中点，并量出 $IJ$
以验证该定理。不测量其他任何长度，
计算 $AIJ$ 的周长。
\end{exercise}

\begin{exercise}[$\star$]\label{exo:g8:midpoints:4}
$IJK$ 是 $ABC$ 的中点三角形，其三边分别为 $10$、
$12$ 和 $16$ cm。求 $IJK$ 的三边及其周长。
两个周长之间有何关系？
\end{exercise}

\begin{exercise}[$\star$]\label{exo:g8:midpoints:5}
一个三角形的面积为 $36$ cm$^2$。其中点三角形的
面积是多少？四个小三角形各自的面积又是多少？
\end{exercise}

\begin{exercise}[$\star\star$]\label{exo:g8:midpoints:6}
在三角形 $ABC$ 中，$I$ 为 $[AB]$ 的中点，$J$ 为
$[AC]$ 的中点，$K$ 为 $[BC]$ 的中点。证明 $IJKB$ 是平行四边形。
（比较 $[IJ]$ 与 $[BK]$：是否平行？是否相等？）
\end{exercise}

\begin{exercise}[$\star\star$]\label{exo:g8:midpoints:7}
一张三角形帆的最下方一条边长 $6.4$ m。一条加固缝线
连接另外两条边的中点，需要多长的缝线？
若改为连接各边上从顶点起算四分之一处的点，
其长度又是多少？（先作猜想并画图，
证明用泰勒斯定理，见\cref{ch:g9:thales}。）
\end{exercise}

\begin{exercise}[$\star\star$]\label{exo:g8:midpoints:8}
在三角形 $ABC$ 中，$A$ 处为直角，$M$ 为
斜边 $[BC]$ 的中点，平行于 $(AB)$ 且过 $M$ 的直线
交 $[AC]$ 于 $N$。
\begin{enumerate}
  \item 证明 $N$ 为 $[AC]$ 的中点。
  \item 说明 $(MN)$ 为何垂直于 $(AC)$。
\end{enumerate}
\end{exercise}

\begin{exercise}[$\star\star\star$]\label{exo:g8:midpoints:9}
设 $ABCD$ 为\emph{任意}四边形，$I$、$J$、$K$、$L$ 分别为
$[AB]$、$[BC]$、$[CD]$、$[DA]$ 的中点（猜想见
\cref{exo:g7:central:11}）。
\begin{enumerate}
  \item 在三角形 $ABC$ 中，将中点定理用于线段
        $[IJ]$，在三角形 $ACD$ 中用于线段 $[LK]$，
        并将二者分别与对角线 $[AC]$ 比较。
  \item 由此得出 $IJKL$ 总是平行四边形。
\end{enumerate}
\end{exercise}

\section{问题：三条中线与重心}

\begin{problem}[{周末问题 --- 三角形的三条中线交于同一点，
该点位于每条中线距顶点三分之二处}]
\label{pb:g8:midpoints:1}
将一个三角形从硬纸板上剪下，它
可以平放在铅笔尖的一个特殊点上而保持平衡，这一点是
三角形的\emph{重心}\index{重心}。
本题求出该点。三角形的一条\emph{中线}\index{中线}是
连接一个顶点与对边\emph{中点}的线段。你将
借助中点定理（两次使用，方法出人意料）证明
三条中线都经过同一点，
并精确地求出该点的位置。

\medskip
\noindent\textbf{第一部分 --- 热身。}
\begin{enumerate}
  \item 画一个较大的三角形 $ABC$（形状不限），
        作 $I$ 为 $[AB]$ 的中点，$J$ 为 $[AC]$ 的中点，$K$ 为
        $[BC]$ 的中点，并画出三条中线 $[AK]$、$[BJ]$、$[CI]$。
        你观察到什么？
  \item 证明：中线将三角形分成两个三角形，
        且二者面积相等。（比较底与高，并利用面积公式：
        底 $\times$ 高 $\div\ 2$。）对于
        面积为 $36$ cm$^2$ 的三角形，两部分的
        面积各是多少？
  \item 在三角形 $ABC$ 中，
        由\cref{thm:g8:midpoints:direct}，关于线段
        $[IJ]$ 可得出什么结论？
\end{enumerate}

\medskip
\noindent\textbf{第二部分 --- 两条中线在距顶点
三分之二处相交。}
中线 $[BJ]$ 与 $[CI]$ 相交于一点，记为 $G$。设
$M$ 为 $[GB]$ 的中点，$N$ 为 $[GC]$ 的中点。
\begin{enumerate}[resume]
  \item 在\emph{三角形 $GBC$}中应用中点定理：
        它关于线段 $[MN]$ 可得出什么结论？
  \item 由此推出 $(IJ) \parallel (MN)$ 且 $IJ = MN$，
        并由\cref{thm:g7:central:recognize}可知
        四边形 $IJNM$ 是平行四边形。
  \item 说明为何 $I$ 与 $N$ 都在中线 $(CI)$ 上，
        $J$ 与 $M$ 都在中线 $(BJ)$ 上 --- 因此
        平行四边形 $IJNM$ 的对角线为 $[IN]$ 与
        $[JM]$，且二者恰好相交于 $G$。由平行四边形的
        \emph{定义}，关于 $G$ 可得出什么结论？
  \item 由此推出 $BM = MG = GJ$，并得出：
        \[
        BG = 2 \times GJ
        \]
        --- 点 $G$ 位于中线 $[BJ]$ 上，距顶点 $B$ 的
        距离为全长的三分之二。请陈述并证明
        在中线 $[CI]$ 上的类似结论。
  \item 现在忽略 $[CI]$，对中线 $[BJ]$ 与 $[AK]$
        这对中线，重复完全相同的论证：它们的交点
        同样在 $[BJ]$ 上距 $B$ 三分之二处。说明为何
        这说明它必为\emph{同一个}点 $G$ --- 并
        得出三条中线都经过 $G$。
  \item 在中线 $[BJ]$ 长为 $9$ cm 的三角形中，
        $G$ 到 $B$ 和到 $J$ 的距离各是多少？
  \item 用一句话说明 $AG = 2 \times GK$ 同样成立。
\end{enumerate}

\medskip
\noindent\textbf{第三部分 --- 六个相等的小块，以及坐标。}
三条中线将三角形分成六个小三角形，
围绕 $G$ 分布。
\begin{enumerate}[resume]
  \item 在三角形 $GBC$ 中，线段 $[GK]$ 是中线。
        由此推出小三角形 $GBK$ 与 $GKC$ 的
        面积相等。
  \item 利用第 2 题，对三角形 $ABK$ 与 $AKC$（二者均由
        中线 $[AK]$ 将 $ABC$ 分割而成），证明三角形 $ABG$ 与 $ACG$
        面积相等。重复此论证，换用另一条中线，
        可得三角形 $ABG$、$BCG$、$CAG$ 的面积
        均为三角形 $ABC$ 面积的三分之一。
  \item 这三个三角形中的每一个都被一段中线一分为二
        （例如 $[GI]$ 是三角形 $ABG$ 的中线 --- 它是从
        哪个顶点引出的？）。由此可得：围绕 $G$ 的六个小块
        面积全部相等，都是三角形面积的
        六分之一。
  \item 在坐标网格上标出 $A(0, 0)$、$B(6, 0)$ 与
        $C(0, 6)$，其中 $K(3, 3)$ 为 $[BC]$ 的中点。利用第二部分的三分之二定理，
        求 $G$ 在中线 $[AK]$ 上的坐标；
        再验证同一点同样位于中线 $[BJ]$ 的三分之二处，
        其中 $J(0, 3)$。
        将 $G$ 的坐标与下列平均数作比较：
        \[
        \frac{0 + 6 + 0}{3},
        \qquad
        \frac{0 + 0 + 6}{3} .
        \]
  \item 平衡实验的结论：利用第 2 题与第 8 题，说明为何
        纸板三角形能平衡在沿任一中线放置的刀口上 ---
        两侧纸板等量 ---
        并说明铅笔尖为何因此必须放在
        $G$ 处，这正是物理学家所称的
        重心。（严格的平衡证明需要大学分册中的
        积分学；此处相等的面积使其
        在今天显得可信。）
\end{enumerate}
\end{problem}
